% GENERATED FILE. Edit canonical lessons, then run npm run generate:books.
% canonical-source-hash: fnv1a64:8582918080f3cff6
% canonical-lessons: ML-C428-varumanam, ML-C428-vilakkayarram, ML-C428-nanayapperuppam, ML-C428-bajarru, ML-C428-sthitivivarakkanakku, ML-R428-review-421-428

\chapter{Money Matters}
\label{ch:ml-money-matters}
\hlchaptermodality{428}

\begin{chapteropening}
\textbf{By the end of this chapter:} \emph{I can say income, rising prices, inflation, a budget and statistics.}

Complete the last lesson of chapter 428: I can say income, rising prices, inflation, a budget and statistics.
\end{chapteropening}

\section[\texorpdfstring{\ml{വരുമാനം} (varumānaṁ)}{വരുമാനം (varumānaṁ)}]{\ml{വരുമാനം} (varumānaṁ) — income}

\label{lesson:ML-C428-varumanam}



\begin{quote}
Before the new one: say the Malayalam for a slum, then the Malayalam for literacy.
\end{quote}

\subsection*{You'll want to know: \ml{വരുമാനം}}
\textbf{\ml{വരുമാനം}} — \emph{varumānaṁ} — \textquotedblleft{}income\textquotedblright{}.

\begin{quote}
\ml{കൃഷിയിൽ} \ml{നിന്നുള്ള} \ml{വരുമാനം} \ml{കുറഞ്ഞു}.
\end{quote}

\begin{quote}
\emph{kṛṣiyil ninnuḷḷa varumānaṁ kuṟaññu.}
\end{quote}

\begin{quote}
The income from farming fell.
\end{quote}

Monthly income is \ml{മാസവരുമാനം}.

\begin{cousinweb}
\textbf{\ml{വരുമാനം}} (\emph{varumānaṁ}) is usually traced to the native verb \ml{വരുക}, \textquotedblleft{}to come\textquotedblright{} (DEDR 5270), whose noun \ml{വരവ്} DEDR also glosses \textquotedblleft{}income\textquotedblright{}: income is what comes in.
\end{cousinweb}

\begin{grammarlens}[title={What you've built}]
One more word for giving reasons and opinions.
\end{grammarlens}

\subsection*{Your turn}
\begin{itemize}
\raggedright
  \item \emph{Say it:} \emph{varumānaṁ}
  \item \emph{Say it:} \emph{varumānaṁ}, once more
  \item \emph{Say it:} \emph{kṛṣiyil ninnuḷḷa varumānaṁ kuṟaññu.} — The income from farming fell.
  \item \emph{Recall:} say the Malayalam for a slum, then the Malayalam for literacy, then say \emph{varumānaṁ} again
\end{itemize}

\begin{tcolorbox}[breakable,colback=teal!4,colframe=teal!35!black,title={Before you move on}]
What does \ml{വരുമാനം} mean? (\textquotedblleft{}Income\textquotedblright{}.) Say it once more.
\end{tcolorbox}
\section[\texorpdfstring{\ml{വിലക്കയറ്റം} (vilakkayaṟṟaṁ)}{വിലക്കയറ്റം (vilakkayaṟṟaṁ)}]{\ml{വിലക്കയറ്റം} (vilakkayaṟṟaṁ) — rising prices}

\label{lesson:ML-C428-vilakkayarram}



\begin{quote}
Before the new one: say the Malayalam for literacy, then the Malayalam for income.
\end{quote}

\subsection*{You'll want to know: \ml{വിലക്കയറ്റം}}
\textbf{\ml{വിലക്കയറ്റം}} — \emph{vilakkayaṟṟaṁ} — \textquotedblleft{}rising prices\textquotedblright{}.

\begin{quote}
\ml{വിലക്കയറ്റം} \ml{കാരണം} \ml{ജീവിതം} \ml{ബുദ്ധിമുട്ടായി}.
\end{quote}

\begin{quote}
\emph{vilakkayaṟṟaṁ kāraṇaṁ jīvitaṁ buddhimuṭṭāyi.}
\end{quote}

\begin{quote}
Life has become hard because of rising prices.
\end{quote}

Its opposite is \ml{വിലക്കുറവ്}, a fall in prices.

\begin{cousinweb}
\textbf{\ml{വിലക്കയറ്റം}} (\emph{vilakkayaṟṟaṁ}) is a native Malayalam compound: \ml{വില}, \textquotedblleft{}price\textquotedblright{} (DEDR 5421), and \ml{കയറ്റം}, \textquotedblleft{}an increase, a climb\textquotedblright{} (DEDR 1255): a climb in prices.
\end{cousinweb}

\begin{grammarlens}[title={What you've built}]
One more word for giving reasons and opinions.
\end{grammarlens}

\subsection*{Your turn}
\begin{itemize}
\raggedright
  \item \emph{Say it:} \emph{vilakkayaṟṟaṁ}
  \item \emph{Say it:} \emph{vilakkayaṟṟaṁ}, once more
  \item \emph{Say it:} \emph{vilakkayaṟṟaṁ kāraṇaṁ jīvitaṁ buddhimuṭṭāyi.} — Life has become hard because of rising prices.
  \item \emph{Recall:} say the Malayalam for literacy, then the Malayalam for income, then say \emph{vilakkayaṟṟaṁ} again
\end{itemize}

\begin{tcolorbox}[breakable,colback=teal!4,colframe=teal!35!black,title={Before you move on}]
What does \ml{വിലക്കയറ്റം} mean? (\textquotedblleft{}Rising prices\textquotedblright{}.) Say it once more.
\end{tcolorbox}
\section[\texorpdfstring{\ml{നാണയപ്പെരുപ്പം} (nāṇayapperuppaṁ)}{നാണയപ്പെരുപ്പം (nāṇayapperuppaṁ)}]{\ml{നാണയപ്പെരുപ്പം} (nāṇayapperuppaṁ) — inflation}

\label{lesson:ML-C428-nanayapperuppam}



\begin{quote}
Before the new one: say the Malayalam for income, then the Malayalam for rising prices.
\end{quote}

\subsection*{You'll want to know: \ml{നാണയപ്പെരുപ്പം}}
\textbf{\ml{നാണയപ്പെരുപ്പം}} — \emph{nāṇayapperuppaṁ} — \textquotedblleft{}inflation\textquotedblright{}.

\begin{quote}
\ml{നാണയപ്പെരുപ്പം} \ml{കാരണം} \ml{പണത്തിന്റെ} \ml{വില} \ml{കുറഞ്ഞു}.
\end{quote}

\begin{quote}
\emph{nāṇayapperuppaṁ kāraṇaṁ paṇattinṟe vila kuṟaññu.}
\end{quote}

\begin{quote}
Because of inflation, the value of money fell.
\end{quote}

A technical word of the news; in talk people say \ml{വിലക്കയറ്റം}.

\begin{cousinweb}
\textbf{\ml{നാണയപ്പെരുപ്പം}} (\emph{nāṇayapperuppaṁ}) is a Malayalam compound: \ml{നാണയം}, \textquotedblleft{}a coin\textquotedblright{}, from Sanskrit \emph{nāṇaka}, \textquotedblleft{}a coin\textquotedblright{} (CDIAL 7050), and \ml{പെരുപ്പം}, \textquotedblleft{}a swelling, an increase\textquotedblright{}, native Dravidian (DEDR 4411): the swelling of money.
\end{cousinweb}

\begin{grammarlens}[title={What you've built}]
One more word for giving reasons and opinions.
\end{grammarlens}

\subsection*{Your turn}
\begin{itemize}
\raggedright
  \item \emph{Say it:} \emph{nāṇayapperuppaṁ}
  \item \emph{Say it:} \emph{nāṇayapperuppaṁ}, once more
  \item \emph{Say it:} \emph{nāṇayapperuppaṁ kāraṇaṁ paṇattinṟe vila kuṟaññu.} — Because of inflation, the value of money fell.
  \item \emph{Recall:} say the Malayalam for income, then the Malayalam for rising prices, then say \emph{nāṇayapperuppaṁ} again
\end{itemize}

\begin{tcolorbox}[breakable,colback=teal!4,colframe=teal!35!black,title={Before you move on}]
What does \ml{നാണയപ്പെരുപ്പം} mean? (\textquotedblleft{}Inflation\textquotedblright{}.) Say it once more.
\end{tcolorbox}
\section[\texorpdfstring{\ml{ബജറ്റ്} (bajaṟṟŭ)}{ബജറ്റ് (bajaṟṟŭ)}]{\ml{ബജറ്റ്} (bajaṟṟŭ) — a budget}

\label{lesson:ML-C428-bajarru}



\begin{quote}
Before the new one: say the Malayalam for rising prices, then the Malayalam for inflation.
\end{quote}

\subsection*{You'll want to know: \ml{ബജറ്റ്}}
\textbf{\ml{ബജറ്റ്}} — \emph{bajaṟṟŭ} — \textquotedblleft{}a budget\textquotedblright{}.

\begin{quote}
\ml{ധനമന്ത്രി} \ml{ഇന്ന്} \ml{ബജറ്റ്} \ml{അവതരിപ്പിച്ചു}.
\end{quote}

\begin{quote}
\emph{dhanamantri innŭ bajaṟṟŭ avatarippiccu.}
\end{quote}

\begin{quote}
The finance minister presented the budget today.
\end{quote}

A family budget is \ml{കുടുംബ} \ml{ബജറ്റ്}.

\begin{cousinweb}
\textbf{\ml{ബജറ്റ്}} (\emph{bajaṟṟŭ}) was borrowed from English \emph{budget}.
\end{cousinweb}

\begin{grammarlens}[title={What you've built}]
One more word for giving reasons and opinions.
\end{grammarlens}

\subsection*{Your turn}
\begin{itemize}
\raggedright
  \item \emph{Say it:} \emph{bajaṟṟŭ}
  \item \emph{Say it:} \emph{bajaṟṟŭ}, once more
  \item \emph{Say it:} \emph{dhanamantri innŭ bajaṟṟŭ avatarippiccu.} — The finance minister presented the budget today.
  \item \emph{Recall:} say the Malayalam for rising prices, then the Malayalam for inflation, then say \emph{bajaṟṟŭ} again
\end{itemize}

\begin{tcolorbox}[breakable,colback=teal!4,colframe=teal!35!black,title={Before you move on}]
What does \ml{ബജറ്റ്} mean? (\textquotedblleft{}A budget\textquotedblright{}.) Say it once more.
\end{tcolorbox}
\section[\texorpdfstring{\ml{സ്ഥിതിവിവരക്കണക്ക്} (sthitivivarakkaṇakkŭ)}{സ്ഥിതിവിവരക്കണക്ക് (sthitivivarakkaṇakkŭ)}]{\ml{സ്ഥിതിവിവരക്കണക്ക്} (sthitivivarakkaṇakkŭ) — statistics}

\label{lesson:ML-C428-sthitivivarakkanakku}



\begin{quote}
Before the new one: say the Malayalam for inflation, then the Malayalam for a budget.
\end{quote}

\subsection*{You'll want to know: \ml{സ്ഥിതിവിവരക്കണക്ക്}}
\textbf{\ml{സ്ഥിതിവിവരക്കണക്ക്}} — \emph{sthitivivarakkaṇakkŭ} — \textquotedblleft{}statistics\textquotedblright{}.

\begin{quote}
\ml{വരുമാനം} \ml{കുറഞ്ഞു} \ml{എന്ന്} \ml{സ്ഥിതിവിവരക്കണക്കുകൾ} \ml{കാണിക്കുന്നു}.
\end{quote}

\begin{quote}
\emph{varumānaṁ kuṟaññu ennŭ sthitivivarakkaṇakkukaḷ kāṇikkunnu.}
\end{quote}

\begin{quote}
The statistics show that income has fallen.
\end{quote}

Often in the plural: \ml{സ്ഥിതിവിവരക്കണക്കുകൾ}.

\begin{cousinweb}
\textbf{\ml{സ്ഥിതിവിവരക്കണക്ക്}} (\emph{sthitivivarakkaṇakkŭ}) is a Malayalam compound: \ml{സ്ഥിതി}, \textquotedblleft{}a state, a condition\textquotedblright{}, from Sanskrit \emph{sthiti} (CDIAL 13769); \ml{വിവരം}, \textquotedblleft{}information\textquotedblright{}, usually traced to Sanskrit \emph{vivaraṇa}, \textquotedblleft{}opening up, explanation\textquotedblright{}, from \emph{vivara}, \textquotedblleft{}an opening\textquotedblright{} (CDIAL 11915); and \ml{കണക്ക്}, \textquotedblleft{}an account\textquotedblright{}, whose origin is uncertain.
\end{cousinweb}

\begin{grammarlens}[title={What you've built}]
One more word for giving reasons and opinions.
\end{grammarlens}

\subsection*{Your turn}
\begin{itemize}
\raggedright
  \item \emph{Say it:} \emph{sthitivivarakkaṇakkŭ}
  \item \emph{Say it:} \emph{sthitivivarakkaṇakkŭ}, once more
  \item \emph{Say it:} \emph{varumānaṁ kuṟaññu ennŭ sthitivivarakkaṇakkukaḷ kāṇikkunnu.} — The statistics show that income has fallen.
  \item \emph{Recall:} say the Malayalam for inflation, then the Malayalam for a budget, then say \emph{sthitivivarakkaṇakkŭ} again
\end{itemize}

\begin{tcolorbox}[breakable,colback=teal!4,colframe=teal!35!black,title={Before you move on}]
What does \ml{സ്ഥിതിവിവരക്കണക്ക്} mean? (\textquotedblleft{}Statistics\textquotedblright{}.) Say it once more.
\end{tcolorbox}
\section[(dialogue)]{Words from chapters 421–428 — a review}

\label{lesson:ML-R428-review-421-428}



\begin{quote}
Say the Malayalam for administration and to govern.
\end{quote}

\subsection*{Your turn}
Say these aloud:

\begin{itemize}
\raggedright
  \item administration, to govern, a minister, a chief minister, a governor
  \item a legislative assembly, a member, a representative, to represent, the opposition
  \item a policy, to implement, to pass a law, to amend, an oath
  \item to reform, an appointment, democracy, authority, a leader
  \item a protest, to protest, a strike, a general strike, a march
  \item to boycott, to demand, a memorandum, to submit, a petition
  \item poverty, a poor man, starvation, a slum, literacy
  \item income, rising prices, inflation, a budget, statistics
\end{itemize}

\begin{tcolorbox}[breakable,colback=teal!4,colframe=teal!35!black,title={Before you move on}]
Administration? (\textbf{\ml{ഭരണം}}, \emph{bharaṇaṁ}.) To govern? (\textbf{\ml{ഭരിക്കുക}}, \emph{bharikkuka}.) A minister? (\textbf{\ml{മന്ത്രി}}, \emph{mantri}.) A chief minister? (\textbf{\ml{മുഖ്യമന്ത്രി}}, \emph{mukhyamantri}.) A governor? (\textbf{\ml{ഗവർണർ}}, \emph{gavarṇar}.) A legislative assembly? (\textbf{\ml{നിയമസഭ}}, \emph{niyamasabha}.) A member? (\textbf{\ml{അംഗം}}, \emph{aṁgaṁ}.) A representative? (\textbf{\ml{പ്രതിനിധി}}, \emph{pratinidhi}.) To represent? (\textbf{\ml{പ്രതിനിധീകരിക്കുക}}, \emph{pratinidhīkarikkuka}.) The opposition? (\textbf{\ml{പ്രതിപക്ഷം}}, \emph{pratipakṣaṁ}.) A policy? (\textbf{\ml{നയം}}, \emph{nayaṁ}.) To implement? (\textbf{\ml{നടപ്പാക്കുക}}, \emph{naṭappākkuka}.) To pass a law? (\textbf{\ml{പാസാക്കുക}}, \emph{pāsākkuka}.) To amend? (\textbf{\ml{ഭേദഗതി} \ml{ചെയ്യുക}}, \emph{bhēdagati ceyyuka}.) An oath? (\textbf{\ml{സത്യപ്രതിജ്ഞ}}, \emph{satyapratijña}.) To reform? (\textbf{\ml{പരിഷ്കരിക്കുക}}, \emph{pariṣkarikkuka}.) An appointment? (\textbf{\ml{നിയമനം}}, \emph{niyamanaṁ}.) Democracy? (\textbf{\ml{ജനാധിപത്യം}}, \emph{janādhipatyaṁ}.) Authority? (\textbf{\ml{അധികാരം}}, \emph{adhikāraṁ}.) A leader? (\textbf{\ml{നേതാവ്}}, \emph{nētāvŭ}.) A protest? (\textbf{\ml{പ്രതിഷേധം}}, \emph{pratiṣēdhaṁ}.) To protest? (\textbf{\ml{പ്രതിഷേധിക്കുക}}, \emph{pratiṣēdhikkuka}.) A strike? (\textbf{\ml{സമരം}}, \emph{samaraṁ}.) A general strike? (\textbf{\ml{ഹർത്താൽ}}, \emph{harttāl}.) A march? (\textbf{\ml{ജാഥ}}, \emph{jātha}.) To boycott? (\textbf{\ml{ബഹിഷ്കരിക്കുക}}, \emph{bahiṣkarikkuka}.) To demand? (\textbf{\ml{ആവശ്യപ്പെടുക}}, \emph{āvaśyappeṭuka}.) A memorandum? (\textbf{\ml{നിവേദനം}}, \emph{nivēdanaṁ}.) To submit? (\textbf{\ml{സമർപ്പിക്കുക}}, \emph{samarppikkuka}.) A petition? (\textbf{\ml{ഹർജി}}, \emph{harji}.) Poverty? (\textbf{\ml{ദാരിദ്ര്യം}}, \emph{dāridryaṁ}.) A poor man? (\textbf{\ml{ദരിദ്രൻ}}, \emph{daridran}.) Starvation? (\textbf{\ml{പട്ടിണി}}, \emph{paṭṭiṇi}.) A slum? (\textbf{\ml{ചേരി}}, \emph{cēri}.) Literacy? (\textbf{\ml{സാക്ഷരത}}, \emph{sākṣarata}.) Income? (\textbf{\ml{വരുമാനം}}, \emph{varumānaṁ}.) Rising prices? (\textbf{\ml{വിലക്കയറ്റം}}, \emph{vilakkayaṟṟaṁ}.) Inflation? (\textbf{\ml{നാണയപ്പെരുപ്പം}}, \emph{nāṇayapperuppaṁ}.) A budget? (\textbf{\ml{ബജറ്റ്}}, \emph{bajaṟṟŭ}.) Statistics? (\textbf{\ml{സ്ഥിതിവിവരക്കണക്ക്}}, \emph{sthitivivarakkaṇakkŭ}.)
\end{tcolorbox}
